%%% FIELD exact-ratio-decoder-proof
We distinguish the structural target \((f^{-1},G,\pi)\), modulo simultaneous relabeling and componentwise diffeomorphisms, from the observable functional \(\mathscr D(P^0,\ldots,P^n)\).  The proof below shows that the latter recovers the former.

Write \(P_i=\operatorname{pa}_G(i)\).  The structural nonredundancy condition used below is precisely \texttt{ass:causal-minimality}:
\[
 b\in P_i\quad\Longrightarrow\quad
 V_i\mathbin{\not\!\perp\!\!\!\perp}_{P^0}V_b
 \mid V_{P_i\setminus\{b\}}.
\tag{1}
\]
Full observational faithfulness, when additionally imposed, implies (1): the length-one path \(b\to i\) has no internal vertex and is not blocked by conditioning on \(P_i\setminus\{b\}\).  Thus \texttt{ass:faithfulness} is used only for the faithful-overlap corollary at the end, whereas the primary decoder and uniqueness argument use (1).

First relabel the latent coordinates by the intervention environments: set \(\bar V_i=V_{\pi(i)}\), pull \(G\), the mechanisms, and the signs back by \(\pi\), and prove the result with \(\pi=\operatorname{id}\).  This is an allowed simultaneous relabeling.  Strict positivity in \texttt{ass:positive-normalized-smooth-mechanisms} makes every ratio below finite and continuous on the closed cube.  By \texttt{ass:shared-diffeomorphic-mixing}, change of variables under the common map \(f\) cancels its Jacobian.  The one-target replacement in \texttt{ass:one-perfect-intervention-per-node} therefore gives, pointwise on \([0,1]^n\),
\[
 R_i(f(v))=\frac{p^i(v)}{p^0(v)}
 =\frac{q_i(v_i)}{p_i(v_i\mid v_{P_i})}
 =:r_i(v),
 \qquad
 L_i(f(v))=\lambda_i(v):=\log r_i(v).
\tag{2}
\]

We next identify the ancestral order.  Put \(A_i=\operatorname{an}_G(i)\cup\{i\}\).  This is an ancestral set containing \(P_i\cup\{i\}\).  If \(j\ne i\) and \(j\notin\operatorname{an}_G(i)\), integrate either the observational factorization or the factorization after intervention \(j\) over \([0,1]^{[n]\setminus A_i}\) in reverse topological order.  At every eliminated vertex the relevant factor is a normalized conditional density, except at \(j\), where it is the normalized density \(q_j\); no eliminated vertex is a parent of a retained vertex because \(A_i\) is ancestral.  Hence
\[
 p^j_{A_i}(v_{A_i})=p^0_{A_i}(v_{A_i})
 =\prod_{a\in A_i}p_a(v_a\mid v_{P_a}).
\tag{3}
\]
The function \(r_i\) is measurable with respect to \(V_{P_i\cup\{i\}}\).  Equation (3) therefore implies
\[
 \operatorname{Law}_{P^j}(R_i)=\operatorname{Law}_{P^0}(R_i),
 \qquad j\notin\operatorname{an}_G(i),\quad j\ne i,
\tag{4}
\]
and hence \(D_{ji}=0\).  Consequently every arrow of \(H_D\) is an ancestral relation of \(G\).  Conversely, \texttt{def:cover-separated-set} puts every cover of the ancestral order into \(H_D\).  Covers generate the whole finite ancestral order: starting with any comparable pair that is not a cover, insert a strictly intermediate vertex; repetition terminates because \([n]\) is finite and produces a saturated chain of covers.  Therefore
\[
 \operatorname{TC}(H_D)\subseteq\operatorname{TC}(G),
 \qquad
 \operatorname{TC}(G)
 =\operatorname{TC}\bigl(\{j\to i:j\lessdot_G i\}\bigr)
 \subseteq\operatorname{TC}(H_D),
\tag{5}
\]
so \(\operatorname{TC}(H_D)=\operatorname{TC}(G)\).  In particular, \(H_D\) is acyclic.

Fix any topological ordering of \(H_D\), and let \(B_i\) be the set of labels before \(i\).  Equality (5) says that this ordering respects every ancestral relation of \(G\); hence
\[
 P_i\subseteq B_i.
\tag{6}
\]
Moreover, \(B_i\) is ancestral: if \(b\in B_i\) and \(a\) is an ancestor of \(b\), then (5) gives a directed \(H_D\)-path from \(a\) to \(b\), so \(a\) precedes \(b\).

For \(b\in B_i\), the function \(\lambda_b\) in (2) depends only on \(v_b\) and \(v_{P_b}\), all of whose parent coordinates precede \(b\).  Thus
\[
 \Lambda_{B_i}:v_{B_i}\longmapsto
 \bigl(\lambda_b(v_b,v_{P_b})\bigr)_{b\in B_i}
\tag{7}
\]
is triangular in the selected order.  By \texttt{ass:fixed-own-derivative-sign},
\[
 s_b\,\partial_{v_b}\lambda_b>0
\tag{8}
\]
throughout the closed cube.  Continuity and compactness make the absolute value in (8) bounded away from zero.  Given the output of (7), one recovers the coordinates successively: with the earlier parents fixed, the \(b\)-th output is a strictly monotone function of \(v_b\).  Hence \(\Lambda_{B_i}\) is injective.  It is continuous, and its inverse on its image is continuous: if image points converge, compactness supplies a convergent subsequence of their preimages, and continuity and injectivity force every such subsequential limit to be the unique preimage of the limiting point.  Consequently
\[
 \sigma(L_{B_i})=\sigma(V_{B_i}).
\tag{9}
\]

The set \(B_i\cup\{i\}\) is ancestral.  Marginalizing the intervention-\(i\) factorization outside that set as in (3) yields
\[
 p^i_{B_i\cup\{i\}}(v_{B_i},v_i)
 =q_i(v_i)\prod_{b\in B_i}p_b(v_b\mid v_{P_b}).
\tag{10}
\]
Thus, under \(P^i\), \(V_i\) has density \(q_i\) and is independent of \(V_{B_i}\).  If \(\ell_{B_i}=\Lambda_{B_i}(v_{B_i})\), equations (2), (9), and (10) supply the following conditional-distribution version on the model support:
\[
 C_i(t\mid\ell_{B_i})
 =\int_0^1
 \mathbf 1\!\left\{\lambda_i(w,v_{P_i})\le t\right\}q_i(w)\,dw.
\tag{11}
\]
The continuous inverse following (7), the strictly monotone inverse furnished by (8), and continuity of \(q_i\) show that (11) is continuous on its model support, including its relative boundary.  It is therefore the continuous version required in the definition of \(C_i\), rather than only an almost-everywhere version.

Evaluate (11) at \(t=L_i(f(v))=\lambda_i(v_i,v_{P_i})\).  Strict monotonicity, the absence of atoms under the density \(q_i\), and normalization give
\[
 U_i(f(v))=
 \begin{cases}
 Q_i(v_i),&s_i=+1,\\[3pt]
 1-Q_i(v_i),&s_i=-1.
 \end{cases}
\tag{12}
\]
Since \(q_i\) is positive and \(C^3\) on the compact interval, both alternatives in (12) are componentwise \(C^2\) diffeomorphisms of the closed interval, with one-sided derivatives at its boundary.  Because (12) constructs \(U_i\) using environment \(i\), it also aligns that environment with the recovered coordinate it targets.

We now prove exact parent pruning using only (1).  Componentwise bimeasurable bijections preserve conditional independence, so (12) implies, for every \(A\subseteq B_i\),
\[
 U_i\mathbin{\perp\!\!\!\perp}_{P^0}U_{B_i\setminus A}\mid U_A
 \quad\Longleftrightarrow\quad
 V_i\mathbin{\perp\!\!\!\perp}_{P^0}V_{B_i\setminus A}\mid V_A.
\tag{13}
\]
Marginalization of the observational factorization outside the ancestral set \(B_i\cup\{i\}\), followed by division by its strictly positive \(B_i\)-marginal, gives
\[
 p^0(v_i\mid v_{B_i})=p_i(v_i\mid v_{P_i}).
\tag{14}
\]
If \(P_i\subseteq A\), the right side of (14) is measurable with respect to \((v_i,v_A)\) and does not depend on \(v_{B_i\setminus A}\); this is the conditional-density factorization for the right side of (13).

Conversely, suppose that the right side of (13) holds and choose an omitted parent \(b\in P_i\setminus A\).  Decompose \(B_i\setminus A\) into \(\{b\}\) and its remaining variables.  The weak-union consequence of the assumed conditional independence is
\[
 V_i\mathbin{\perp\!\!\!\perp}V_b
 \mid V_{B_i\setminus\{b\}}.
\tag{15}
\]
Put \(Z=P_i\setminus\{b\}\) and \(W=B_i\setminus P_i\).  Equation (15) is \(V_i\mathbin{\perp\!\!\!\perp}V_b\mid(V_Z,V_W)\).  Equation (14), on the other hand, gives the ordered local-Markov relation
\[
 V_i\mathbin{\perp\!\!\!\perp}V_W\mid(V_Z,V_b).
\tag{16}
\]
Strict positivity proves the required intersection step directly.  Indeed, conditional independence first gives the following identities almost everywhere; positivity and continuity of the marginal and conditional densities extend them to every point of the product support:
\[
 p^0(v_i\mid v_b,v_W,v_Z)=p^0(v_i\mid v_W,v_Z),
 \qquad
 p^0(v_i\mid v_b,v_W,v_Z)=p^0(v_i\mid v_b,v_Z).
\tag{17}
\]
For each fixed \((v_i,v_Z)\), the left member of the resulting equality \(p^0(v_i\mid v_W,v_Z)=p^0(v_i\mid v_b,v_Z)\) depends only on \(v_W\), while the right member depends only on \(v_b\).  Full positive product support permits arbitrary \((v_b,v_W)\), so both sides are constant in those arguments.  Hence
\[
 V_i\mathbin{\perp\!\!\!\perp}V_b\mid V_Z,
\tag{18}
\]
contradicting causal minimality (1).  Therefore no parent can be omitted.  Together with the sufficient direction this proves
\[
 U_i\mathbin{\perp\!\!\!\perp}_{P^0}U_{B_i\setminus A}\mid U_A
 \quad\Longleftrightarrow\quad P_i\subseteq A.
\tag{19}
\]
The admissible subsets are exactly the supersets of \(P_i\); consequently, \(P_i\) is their unique inclusion-minimal member.  This proves that \texttt{def:population-decoder} returns \(G\), and the argument applies to every topological ordering of \(H_D\).

It remains to prove uniqueness among compatible representations; cover separation is not imposed on the competitor.  Relabel both representations by their target bijections, write the competing DAG as \(\widetilde G\), its coordinates as \(z\), its mixing map as \(\widetilde f\), and its ratios and intervention distribution functions with tildes.  The implication (4) holds for the competitor because it used only factorization, normalization, positivity, and a perfect intervention.  Therefore every observable arrow in \(H_D\) is also an ancestral relation in \(\widetilde G\), and (5) gives
\[
 \operatorname{TC}(G)=\operatorname{TC}(H_D)
 \subseteq\operatorname{TC}(\widetilde G).
\tag{20}
\]
Choose a topological order of \(\widetilde G\).  By (20), it is also a topological order of \(G\) and of \(H_D\); use this common order below, writing \(j>i\) when \(j\) is later than \(i\).

The transition map
\[
 T=f^{-1}\circ\widetilde f:[0,1]^n\longrightarrow[0,1]^n
\tag{21}
\]
is a \(C^2\) diffeomorphism.  Equality of the observed laws makes their likelihood ratios equal almost everywhere.  Their latent versions are continuous, and the positive observational density assigns positive probability to every nonempty relatively open set, so the equality extends everywhere:
\[
 r_i(T(z))=\widetilde r_i(z),\qquad i\in[n].
\tag{22}
\]
We prove inductively in the common order that \(T_i\) depends only on \(z_1,\ldots,z_i\).  If \(j>i\), the right side of (22) has zero \(z_j\)-derivative because it depends only on \(z_i\) and the earlier \(\widetilde G\)-parents of \(i\).  The chain rule and the induction hypothesis for every \(a\in P_i\) give
\[
 \begin{aligned}
 0=\partial_{z_j}\widetilde r_i(z)
 &=\partial_{v_i}r_i(T(z))\,\partial_{z_j}T_i(z)
 +\sum_{a\in P_i}\partial_{v_a}r_i(T(z))\,\partial_{z_j}T_a(z)\\
 &=\partial_{v_i}r_i(T(z))\,\partial_{z_j}T_i(z).
 \end{aligned}
\tag{23}
\]
By (8), \(\partial_{v_i}r_i=r_i\partial_{v_i}\lambda_i\ne0\).  Hence \(\partial_{z_j}T_i=0\), proving that \(T\) is lower triangular.  Since \(DT\) is nonsingular and triangular, every diagonal derivative \(\partial_{z_i}T_i\) is nonzero; its sign is constant on the connected cube.

Every prefix map \(T_{\le i}\) is a diffeomorphism of the corresponding prefix cube.  Indeed, strict monotonicity in each diagonal coordinate gives injectivity successively.  For surjectivity, fix a desired output prefix, extend it arbitrarily to a point of \([0,1]^n\), use surjectivity of \(T\), and then use triangularity to see that the preimage's input prefix maps to the desired prefix.  Compare prefix densities in environment \(i\).  Normalizing all later factors in the common topological order gives
\[
 p^i_{\le i}(v_{\le i})=p^0_{<i}(v_{<i})q_i(v_i),
 \qquad
 \widetilde p^i_{\le i}(z_{\le i})
 =\widetilde p^0_{<i}(z_{<i})\widetilde q_i(z_i).
\tag{24}
\]
Triangularity, prefix bijectivity, and equality of the observational prefix laws give
\[
 \widetilde p^0_{<i}(z_{<i})
 =p^0_{<i}(T_{<i}(z_{<i}))\,|\det DT_{<i}(z_{<i})|.
\tag{25}
\]
Apply change of variables under \(T_{\le i}\) to the two environment-\(i\) identities in (24), substitute (25), and cancel the strictly positive common factor.  The result is the scalar fiber identity
\[
 q_i(T_i(z_{<i},z_i))
 \left|\partial_{z_i}T_i(z_{<i},z_i)\right|
 =\widetilde q_i(z_i).
\tag{26}
\]

Every fiber \(z_i\mapsto T_i(z_{<i},z_i)\) maps \([0,1]\) onto \([0,1]\).  Indeed, fix the input prefix \(z_{<i}\) and a desired output \(v_i\); extend \((T_{<i}(z_{<i}),v_i)\) to an output vector and invoke surjectivity of \(T\).  Prefix injectivity forces the resulting preimage to have prefix \(z_{<i}\).  Thus an increasing fiber has endpoints \(0,1\), and a decreasing fiber has endpoints \(1,0\).  Integrating (26) from \(0\) to \(z_i\) now yields
\[
 Q_i(T_i(z_{<i},z_i))=
 \begin{cases}
 \widetilde Q_i(z_i),&\partial_{z_i}T_i>0,\\[3pt]
 1-\widetilde Q_i(z_i),&\partial_{z_i}T_i<0.
 \end{cases}
\tag{27}
\]
Because \(Q_i\) is strictly increasing, (27) makes \(T_i\) a function of \(z_i\) alone.  Positivity and smoothness of the two intervention densities make this scalar map a \(C^2\) diffeomorphism.  Induction proves that \(T\) is componentwise.

It remains only to compare the DAGs.  Apply the already-proved equivalence (19) to \(G\) and to \(\widetilde G\) in the common order.  Componentwise bimeasurable transformations preserve all conditional independences used by the pruning functional.  Thus the family of admissible subsets \(A\subseteq B_i\) is the same for both representations.  Causal minimality makes its unique inclusion-minimal member equal to \(P_i\) in the first representation and to \(\operatorname{pa}_{\widetilde G}(i)\) in the second.  Therefore
\[
 \operatorname{pa}_G(i)=\operatorname{pa}_{\widetilde G}(i)
 \quad\text{for every }i,
\tag{28}
\]
so \(G=\widetilde G\).  Undoing the two target relabelings proves uniqueness up to simultaneous relabeling and componentwise \(C^2\) diffeomorphisms.  Restoring the original \(\pi\) in (5), (6), (12), and (19) gives every displayed identification conclusion in the statement and identifies the target assignment through the environment-to-coordinate alignment.

We finally verify the comparison clauses.  The established scope records \texttt{lem:von-kugelgen-comparator-scope}, \texttt{lem:cauca-comparator-scope}, and \texttt{lem:yao-comparator-scope} supply the published intervention-count and supplied-structure facts.  The present theorem assumes positive \(C^3\) mechanisms on a compact cube and a shared \(C^2\) mixing map, so its conclusion does not extend by this proof to the anchor's full \(\mathbb R^n\), \(C^1\) class.  On the overlapping faithful regime, faithfulness implies (1), and the theorem uses one intervention per node rather than the paired design in the arbitrary-dimensional comparator.  Within the present regime, the proof of (19) and (28) uses only causal minimality, which establishes the stated structural extension without describing the causal-minimal class as a subclass or specialization of the anchor's faithful class.

For the bivariate witness clause, let \(\nu=\operatorname{Law}_{P^0}(R_i)-\operatorname{Law}_{P^j}(R_i)\) and define
\[
 \varphi(r)=\int_{\mathbb R}k(r,t)\,d\nu(t).
\tag{29}
\]
Boundedness and continuity of the Gaussian kernel make \(\varphi\) bounded and continuous by dominated convergence.  The defining double-integral formula for MMD and Fubini's theorem, applicable because \(|k|\le1\), give
\[
 \mathbb E_{P^0}\varphi(R_i)-\mathbb E_{P^j}\varphi(R_i)
 =\int_{\mathbb R}\!\int_{\mathbb R}k(r,t)\,d\nu(t)\,d\nu(r)
 =D_{ji}^{2}>0.
\tag{30}
\]
Thus Gaussian-MMD positivity supplies the continuous law-separation witness used in the bivariate comparator.  Together with the established scope records, this proves every comparison clause.  Finally, the implication from full observational faithfulness to (1), proved at the start, makes the faithful version an immediate corollary of the causal-minimal decoder.
